%=====================================================================
\section{Translation}
%=====================================================================

Translation is performed by the translation model (Table~\ref{tab:models}). The model receives a one-line instruction, \emph{``Translate the following text into English''}, followed by the Bengali text. The translation is generated greedily, without sampling, so the same input always gives the same output.

Two versions of the sentence are translated:
\begin{itemize}[nosep]
  \item The \textbf{clean} sentence is the transcript without any markers.
  \item The \textbf{marked} sentence contains the pause markers from the alignment and prosody marking stage.
\end{itemize}
Comparing the two versions shows the effect of the markers on the translation of the example. The marked translation is used for synthesis.

\begin{example}[: translation]
The diagram aligns the marked Bengali with both English translations, and the punctuation and capital letters that the translator added are highlighted. \textbf{Without the marker, the question was no longer marked as a question.} For this clip, the output with the marker is consistent with the intended effect of the pause markers.

\begin{center}
% Worked example, Stage 4: the marked Bengali, split at its long-pause bar, mapped to the
% two sentences of the marked English translation, with the unmarked translation below.
% Punctuation and capital letters that the translator added are shown in red.
\begin{tikzpicture}[
    bengali/.style={draw=blue!55!black, fill=blue!5, rounded corners=3pt, align=center,
                    anchor=north, inner sep=5pt},
    english/.style={draw=green!40!black, fill=green!5, rounded corners=3pt, align=center,
                    anchor=north, inner sep=5pt, font=\small,
                    execute at begin node={\hyphenpenalty 10000\relax}},
    bar/.style={text=red!75!black, font=\Large\bfseries},
    row title/.style={font=\footnotesize\bfseries, anchor=south west},
    maps to/.style={-{Stealth[length=2mm]}, black!60, line width=0.8pt},
    note/.style={font=\scriptsize\itshape, text=violet, align=center},
  ]
  % added punctuation and capitals: bold red letters
  \def\hl#1{\textbf{\color{red!80!black}#1}}

  % Marked Bengali: two phrases separated by one long-pause bar
  \node[row title] at (0.3,0.1) {Marked Bengali};
  \node[bengali, text width=7.4cm] (b1) at (4.2,0) {\bn গ্রেগরির কাছ থেকে এখনও উত্তর পেলাম না};
  \node[bar] at (8.35,0 |- b1.center) {|};
  \node[bengali, text width=5.2cm] (b2) at (11.45,0) {\bn সে কি তাহলে ইংল্যাণ্ডে নেই};

  % Marked English: the bar became a sentence break and was not copied
  \coordinate (row2) at ([yshift=-14mm]b1.south);
  \node[english, text width=7.4cm] (e1) at (4.2,0 |- row2)
    {\hl{I} haven't received a reply from \hl{G}regory yet\hl{.}};
  \node[english, text width=5.2cm] (e2) at (11.45,0 |- row2)
    {\hl{I}s he not in \hl{E}ngland\hl{?}};
  \node[row title] at ([yshift=1mm]0.3,0 |- e1.north) {Marked English};
  \node[note] at ($(8.35,0 |- b1.south)!0.5!(8.35,0 |- e1.north)$)
    {the bar was not copied;\\a sentence break appeared};
  \draw[maps to] (b1.south) -- ([yshift=1mm]e1.north);
  \draw[maps to] (b2.south) -- ([yshift=1mm]e2.north);

  % The same sentence translated without markers
  \node[row title, anchor=north west] (clean) at ([yshift=-5mm]0,0 |- e1.south) {English without markers};
  \node[draw=black!50, fill=black!4, rounded corners=3pt, text width=13.6cm, align=center,
        anchor=north west, inner sep=5pt, font=\small] (c) at ([xshift=0.3cm]clean.south west)
    {i haven't received a reply from gregory yet is he not in england};
  \node[note, text=black!60, anchor=north] at (c.south) {one run-on string: no punctuation and no capital letters};
\end{tikzpicture}

\end{center}
\end{example}
